\section{Corrections and narrowed claims}
\label{sec:corrections}

Answering the reviews uncovered statements of the main paper that are imprecise, unsupported or wrong. We list them with the place where each is treated, so that the paper and this document can be read together without ambiguity. The numbers of Tables~2, 4 and~5 are reproduced exactly and do not change; one row of Table~3 changes; and several conclusions drawn from these tables have to be narrowed.

\begin{enumerate}[leftmargin=*,itemsep=2pt,topsep=3pt]
\item \textbf{Definition of the success rate (Sec.~5, ``Metrics'').} SR in Tables~2--4 thresholds the range error $d$, not the three-dimensional distance to the goal as the text says. Under the endpoint definition the refined stage succeeds on \SiiCamSRe\% of the pairs in the camera frame and \SiiEnuSRe\% in ENU, instead of 83.3\%; the coarse stage on \SiCamSRe\% and \SiEnuSRe\% (\cref{sec:whichdef}, \cref{tab:main_metrics}). Both definitions should be reported, and the headline number of the abstract should be accompanied by the endpoint success.
\item \textbf{The $d$ column} is the error of the predicted distance, not of the displacement vector (\cref{sec:whichdef}).
\item \textbf{Comparison with existing estimators (abstract, Sec.~1, Tables~2 and~4).} The eight baselines are navigation policies, world models and visual representations adapted to regress a relative pose, several of them with a linear probe or a single epoch (\cref{sec:baselines}); they are not ``widely used relative-pose regressors''. Against correspondence-based estimation the margin is small or negative: a SIFT pipeline with an essential matrix and the same flight-height prior reaches \sfSiftSRd\% / \sfSiftSRe\% (\SRd{} / \SRe) on the test pairs of Table~2, between our coarse stage (\sfSiSRd\% / \sfSiSRe\%) and our refined stage (\sfSiiSRd\% / \sfSiiSRe\%); RoMa with a rotation search is level with the refined stage on the pairs on which it was run; and on the real photographs both are more accurate than DreamNav (\grSiftSRe\% of endpoint success for the SIFT pipeline against \grSiiSRe\%; \cref{sec:geometric,sec:real}). These rows belong in Tables~2 and~4, and the claim of superiority should be restricted to the methods of Table~2 and to the refined stage on PairUAV.
\item \textbf{Table~3, first row.} ``Frozen SuperGlue'' is a regression from pooled MASt3R statistics and uses no correspondences (\cref{sec:inventory}).
\item \textbf{Table~3, third row.} ``+ absolute heading'' should read $d{=}$\fOneD\,m and SR${=}$\fOneSRd\%, not 2.25\,m and 35.1\%. Consequently the heading does not reduce the distance error, the zoom re-matching adds 8.0 points and not 38.0, and the corresponding sentences of Sec.~5.2 are withdrawn (\cref{sec:inventory}).
\item \textbf{Role of the plane-induced warp (title, abstract, Secs.~1 and~4.3).} The paper presents the ground-plane homography as what makes the refinement work. A compute-matched second matching pass that aligns the images by the coarse rotation only recovers \ctRotShareD\% of the gain in \SRd{} and \ctRotShareE\% of the gain in endpoint success; the plane-induced part adds \ctPlaneGainSRe{} points of endpoint success \ci{\ctPlaneGainSReLo}{\ctPlaneGainSReHi}, and the verification scores add nothing measurable. The contribution of the refinement stage should be described as pose-guided re-matching at the original resolution, of which the plane-induced warp is one, modestly better, form of alignment (\cref{sec:controls}).
\item \textbf{Abstract, internal \vs physical loop.} The increase from 38.8\% to 81.9\% is obtained with 2.2 \emph{physical} steps in simulation, each followed by a new rendered observation. The sentences ``Iterating the refinement 2.2 times [...] the loop closes inside the estimator rather than in the air'' conflate this with the internal refinement and are withdrawn. Repeating the refinement on a fixed pair, without moving, changes \SRe{} by \itThreeASRe{} \ci{\itThreeASReLo}{\itThreeASReHi} to \itThreeBSRe{} \ci{\itThreeBSReLo}{\itThreeBSReHi} points after three rounds, depending on the protocol (\cref{sec:iteration,sec:internalphysical}).
\item \textbf{Table~5, ``Ours (Stage~II)''.} This row uses, below 5\,m of estimated remaining distance, the ground-plane translation of \cref{sec:shortrange}, which the main paper does not describe. Averaged over \clSeeds{} draws of the actuation noise, the refined stage without it reaches \clSiiPlainSR\% at $\tau{=}1$\,m, against \clSiSR\% for the coarse stage and \clSiiSR\% with it: the gain of Table~5 comes from this solver and not from the refinement (\cref{sec:cluncertainty}). A classical image-based servo reaches \svIbvsFinal\% in the same simulator when it is given 20 steps (\cref{sec:servo}).
\item \textbf{Table~4, flight height.} The two DreamNav rows scale the predicted displacement by the ratio of the flight height reported by the aircraft (60\,m) to the height of the training data; the baselines receive no such input. Without the flight height the refined stage has a median position error of \rlNoHeightPos\,m (\cref{sec:real}).
\item \textbf{Sec.~5.3, ``generalizes well''.} This holds for the pairs of Table~4, whose baseline-to-height ratio lies within the training range, with the flight height given and in the camera frame. Beyond that range the endpoint success is \rlSiiOutSRe\%, and with the learned absolute heading it is \rlEnuSRe\% (\cref{sec:real}).
\item \textbf{``Unseen scenes'' and the ENU results.} The test pairs come from held-out \emph{buildings}, whose surroundings are not held out: the median test building stands \hdDistMedian\,m from a training building. The geometric part of the method is insensitive to this, but the learned absolute heading is not. On six sites removed entirely from the training of the heading network, the ENU endpoint success falls from \hdSdSeenSRe\% to \hdSdUnseenSRe\%, while the camera-frame success on the same pairs is \hdSdCamSRe\% (\cref{sec:headsite}); on the real photographs the heading prior is at $\hdRealFusedMed^\circ$ of median error (\cref{sec:headsens}). ``Unseen scenes'' should read ``unseen buildings'', the camera-frame results should be the primary ones, and the $\Delta x$, $\Delta y$ columns of Tables~2 and~3 should be presented as conditional on a heading prior that is specific to the sites of PairUAV. \todoauthors{Sec.~5 describes the split of the complete dataset, whereas the experiments of the paper and of this document use \nTest{} test pairs from \nBuildings{} held-out buildings and training pairs from 1{,}122 other buildings; please reconcile the two descriptions.}
\item \textbf{Sec.~4.2, absolute heading.} The refined output combines two heading networks, not one (\cref{sec:absh}).
\item \textbf{Eq.~(1).} The prior heading change in the last term lacks its bar (\cref{eq:map}).
\item \textbf{Figure~1 and the conclusion.} The experiments establish pose estimation from views taken about 220\,m above the ground and a simulated final approach without obstacles, dynamics or a flight controller. They do not establish close-range flight along a facade (\cref{sec:planarity}), obstacle avoidance, or deployment on an aircraft (\cref{sec:clscope}); the last sentence of the abstract and of the conclusion should be read as motivation.
\end{enumerate}

\noindent\textbf{Dependences that the paper should state.} (i) Metric scale relies on the regularity of the flight height in PairUAV, or on the flight height being given; this holds for our system and for every geometric baseline (\cref{sec:inventory,sec:georesults,sec:real}). (ii) The ENU output relies on an absolute heading learned from image appearance, which is accurate only near the training buildings and under the imaging conditions of the rendered imagery; the camera-frame output does not (\cref{sec:heading}). (iii) The calibration $K$, zero roll and tilts within $0^\circ$--$45^\circ$ are assumed (\cref{sec:conv,sec:coarse}). (iv) The pairs overlap: an estimate is returned for every pair, also for those without common content (\cref{sec:unfiltered}).

\noindent\textbf{What remains.} After these corrections the supported contributions are: the five-degree-of-freedom benchmark with metric labels; a pose-guided second matching pass at the original resolution that raises the endpoint success of a learned sparse matcher by \gainSReCam{} points and brings its translation direction to the level of a high-resolution classical pipeline; and a learned fusion of scale cues that estimates the distance better than a single flight-height cue on PairUAV (\cref{sec:georesults}). The same fusion does not transfer this advantage to the real photographs.
